\documentclass[11pt]{article}
\usepackage[T1]{fontenc}
\usepackage{lmodern}
\usepackage{microtype}
\usepackage[margin=1in]{geometry}
\usepackage{amsmath,amssymb}
\usepackage{xcolor}
\usepackage[hidelinks]{hyperref}

\definecolor{statusblue}{RGB}{29,78,121}
\definecolor{softgray}{RGB}{245,247,249}

\title{A Self-Adjoint Operator Whose Norm Is Not Attained}
\author{Literature-implied status note for arXiv:2004.05496}
\date{9 August 2026}

\begin{document}
\maketitle

\noindent\colorbox{softgray}{\parbox{0.94\linewidth}{
\textbf{Status.} \textcolor{statusblue}{Literature-implied answer, full for
Open Question 1.} The intended norm-attainment question has an affirmative
answer already implicit in arXiv:1209.1218. This is not an original pipeline
result.
}}

\section*{Source question}
N. B. Okelo, \emph{Various notions of norm-attainability in normed spaces},
arXiv:2004.05496, Section 7, p.~12, asks \cite{Okelo}:
\begin{quote}
Let $H$ be a reflexive, dense, separable, infinite dimensional complex
Hilbert space. Does there exist a bounded self-adjoint operator
$A:H\to H$ such that $\|Ax_0\|<\|A\|$ for all $x_0\in H$?
\end{quote}
Definition 2.5 of the same paper (p.~2) defines norm attainment using a
\emph{unit} vector. Thus the mathematically intended question is whether a
bounded self-adjoint operator can satisfy
\[
  \|Ax\|<\|A\| \qquad (\|x\|=1).
\]

\section*{Identification}
S. Shkarin, \emph{Norm attaining operators and pseudospectrum},
arXiv:1209.1218, p.~2, explicitly considers the diagonal operator $D$ on
$\ell_2$ with diagonal entries
\[
  d_n=1-2^{-n}, \qquad n\in\mathbb N,
\]
and states that $D$ has norm $1$ which is not attained \cite{Shkarin}.
This is exactly the requested example. Indeed, $0<d_n<1$ and the diagonal is
real, so $D$ is bounded, positive, and self-adjoint. Moreover,
\[
 \|D\|=\sup_n d_n=1,
\]
whereas for every unit vector $x=(x_n)$,
\[
 \|Dx\|^2=\sum_{n=1}^{\infty}d_n^2|x_n|^2
 <\sum_{n=1}^{\infty}|x_n|^2=1.
\]
The strict inequality holds because at least one $x_n$ is nonzero and every
$d_n^2$ is strictly below $1$. Every separable infinite-dimensional complex
Hilbert space is unitarily isomorphic to $\ell_2$, so conjugating $D$ by such
a unitary gives the example on any $H$ in the question.

\section*{Literal wording}
Without the unit-vector restriction, the displayed quantifier in the source
question is impossible. If $A\ne0$, choose $u$ with $Au\ne0$ and then choose
$t>\|A\|/\|Au\|$; this gives $\|A(tu)\|>\|A\|$. If $A=0$, the requested strict
inequality is $0<0$. Thus the literal reading has a negative answer, while the
intended norm-attainment reading has the affirmative answer above.

\section*{Provenance, search evidence, and scope}
Shkarin's 2012 paper predates Okelo's 2020 question and does not identify
itself as answering that later question. The relation is therefore an
agent-identified implication and belongs in the literature-implied category.
An exact-phrase web search for Okelo's question found the source paper but no
later paper explicitly advertising an answer; the run's local full-text index
then located Shkarin's earlier diagonal example. The example is stated on
p.~2 of the supporting PDF and the source question is on p.~12 of the source
PDF.

This note fully answers Open Question 1 under both possible readings. It does
not address Open Question 2 about when $p$-norm-attainability and
$p$-normality coincide in general Banach spaces.

\textbf{Human-review recommendation.} Confirm that Okelo's Definition 2.5
supplies the unit-sphere convention and that Shkarin's diagonal operator is
being considered over the complex field. No computational dependency is
present.

\begin{thebibliography}{9}
\bibitem{Okelo}
N. B. Okelo,
\emph{Various notions of norm-attainability in normed spaces},
arXiv:2004.05496 (2020), Definition 2.5 and Section 7.

\bibitem{Shkarin}
S. Shkarin,
\emph{Norm attaining operators and pseudospectrum},
arXiv:1209.1218 (2012), p.~2.
\end{thebibliography}

\end{document}
